% !TeX root = ../main.tex
% !TeX program = xelatex
% !TeX encoding = UTF-8
\section{多目标规划与物理反馈求解}

\subsection{费用与系统碳的容差端点及约束扫描}

技术、管网和运行模式选择含整数变量，能流、容量与热状态为连续变量，装配后形成有限规模MILP。费用目标为
\begin{equation}
F_C=C_0+C_{\mathrm{carbon}}+C_{\mathrm{pref}},\qquad F_E=\boldsymbol a_E^\mathsf T\boldsymbol x+E_0,
\label{eq:scoring-objectives}
\end{equation}
其中偏好与惩罚单列，系统碳采用冷热服务及共享电力的一致台账，常数项$E_0$保留。所有方案沿用相同的可靠性、舒适、设备和资源硬边界，防止通过减少服务获得虚假减碳。

为处理单目标的同值替代解，先以费用为主、系统碳为次求一个端点，再交换次序、新建模型求另一端点\supcite{5}。第一阶段获得主目标可行值$v_1$后，第二阶段加入
\begin{equation}
f_1\leq v_1+\delta_{\mathrm{abs}}+\delta_{\mathrm{rel}}|v_1|,
\qquad\min f_2.
\label{eq:scoring-endpoint}
\end{equation}
接口默认$\delta_{\mathrm{abs}}=10^{-7}$，量纲随主目标，$\delta_{\mathrm{rel}}=10^{-8}$；两阶段分别登记状态、时限、原始界与gap。若第一阶段仅返回限时可行解，第二阶段以该可行值为锚点，其身份仍为容差边界内的次目标改进。

随后在端点碳区间内按登记档位求解$\min F_C$、$F_E\leq\varepsilon_E^{\mathrm{eff}}$，其中$\varepsilon_E^{\mathrm{eff}}=\varepsilon_E+10^{-7}+10^{-8}|\varepsilon_E|$。专用扫描入口将碳价置零，赛事入口保留配置中的碳价与偏好项，两类目标分别标注。每档复核碳帽与数值可行性，对可行样本去重并剔除支配点，形成有限采样下的非支配方案集合。碳帽具有明确物理单位，便于直接回答“给定排放额度下能以多少费用满足服务”。

\subsection{实际日选取与硬能量约束下的权重校正}

代表日首先将同源气象、负荷等逐时特征按日展开、标准化，以K-means聚类，再取各簇距中心最近的实际日。随后保留系统及各服务负荷峰值、最大真实相邻小时爬坡、最低气温和已配置的资源风险日；跨午夜最大爬坡保留两侧日期，新增极值日的频数从原簇扣除，使总日数守恒。

频数权重保证覆盖次数，却可能使压缩后的总能量和负荷持续分布偏离原序列。为此，在固定选日集合上仅校正权重。设保留观测含$D$个完整日，每日$H=24$小时，实际日轮廓为$q_{jh}$，$E_j=\sum_hq_{jh}\Delta t$；在观测与代表值联合支撑点$s_i$上计算经验分布$F_o$及加权分布$F_w$，求解
\begin{equation}
\begin{aligned}
\min_{\boldsymbol w}\;&\frac{1}{\bar q}\sum_i(s_{i+1}-s_i)
       |F_w(s_i)-F_o(s_i)|,\\
\text{s.t.}\;&\sum_jw_j=D,\quad \sum_jw_jE_j=E_o,\quad w_j\geq\ell_j,\\
F_w(s_i)=\;&\frac{1}{HD}\sum_jw_j\sum_{h=1}^{H}\mathbf1(q_{jh}\leq s_i).
\end{aligned}
\label{eq:scoring-hard-w1}
\end{equation}
其中$\bar q>0$，$E_o$为同一保留序列的总能量。绝对值用辅助变量线性化，由HiGHS求解标准LP，并独立重算能量、日数、下界和目标残差。冻结六楼规程对极值日取$\ell_j=1$日，其余取$10^{-6}$日；校正不改原观测或日轮廓。

候选还须通过预先登记的统计准入：相对能量误差不超过0.1\%，同口径归一$W_1$不劣于频数基线，保留序列峰值、真实相邻爬坡及原始观测峰值保留率均达到1。失败时保存原因并回退频数权重，回退方案另行验收。事件覆盖由选日与峰值检查保证，权重校正优化既定样本的分布；状态可行性由连续回放检查。

\subsection{固定投资原时序回放与管网反馈}

规划结果保存技术、容量、管网、储热和接入选择，并通过冻结审计固定全部投资；原时序仅重新调度。回放先核对输入指纹、真实时间戳、容量与服务边界，再连续递推建筑及储热状态，独立复算舒适、能量和初末条件。接口覆盖所给连续窗口，三城共同2880小时窗口及其控制策略试验分别登记。统计压缩承担规划计算，原时序审核承担运行准入。

管网反馈进一步处理线性调度解与瞬时热水输送条件的差异。当前实现限定于独立泵星型支路、定物性水和有限个固定流量、回水温度稳态模式。每个模式计算两程散热、远端供水、摩阻和泵电力。令$z_{mt}$表示模式启用，则
\begin{equation}
0\leq q_{mt}\leq Q_m^{\mathrm{in}}z_{mt},\quad
\sum_{m\in b}z_{mt}\leq1;
\qquad z_{mt}=0\quad(m\in\mathcal B^{(k)},\ \forall t),
\label{eq:scoring-network-cut}
\end{equation}
其中$\mathcal B^{(k)}$为第$k$轮发现的不合格模式。小时平均交付量按占空比计算，但流速、压力、泵功率和供回温边界按启用瞬时值核查；源容量与电力接入亦按同时启用模式约束。每轮重新计算源热、交付热、散热及泵电台账，失败模式全时段禁用后重新求解；无新增有效排除、求解失败或轮数耗尽均停止并留痕。有限目录逐步缩减，因此每次有效反馈至少消去一个模式。温差软惩罚单独核算，其作用须由开关对照确认；当前试验的可行性改善来自硬校核排除。该反馈不涵盖输送时滞、启停瞬态或共泵水力平衡。

\subsection{同目标收敛证据与模块衔接}

核心模型经SciPy的\path{milp}接口调用HiGHS\supcite{6,7}。结果同时保存终止状态、可行解、矩阵及整数残差、目标重算和gap。最小化问题的可行目标值$U$与求解器下界$L$采用相同目标，含常数项时共同加回；限时可行解与容差内最优解分别标注。固定设备热泵LP另以精确有理运算处理已编码binary64系数，形成电量费目标的弱对偶下界。该证书与数值解的可行性容差共同说明解的质量，便于不同边界与时限下复核收敛结果。

模块通过输入身份、服务合同和固定资产记录衔接。多目标端点与碳帽扫描形成取舍集合，$W_1$校正提高输入分布保真，连续回放及物理反馈检查服务可达性。第七章以伪代码展开各入口的调用、停止和复现逻辑。
